\chapter{算术理论中的推导}\label{deriv:chap}

在我们证明所有一般递归函数都在~$\Th{Q}$ 中可表示，以及证明不完全性定理时，我们曾断言各种结论在 $\Th{Q}$ 和~$\Th{PA}$ 中是可证的。然而，这些断言的证明只是非形式地给出了论证，并未展示自然演绎中的实际推导。我们将在本章提供其中一些推导。

例如，在 \olref[inc][req][bre]{lem:q-proves-add} 中，我们通过对 $m$ 归纳，证明了对于所有 $n$ 和 $m \in \Nat$，$\Th{Q} \Proves (\num{n} +
\num{m}) = \num{n+m}$。

\begin{proof}[{\olref[inc][req][bre]{lem:q-proves-add}} 的证明]
基始：$m = 0$。此时需要证明的是，对所有 $n$，$\Th{Q} \Proves \num{n} + \num{0} = \num{n+0}$。由于 $\num{0}$ 就是 $\Obj 0$ 而 $\num{n+0}$ 就是 $\num{n}$，这就相当于要证明$\Th{Q} \Proves (\num{n} + \Obj 0) = \num{n}$。推导
\begin{prooftree}
  \AxiomC{$\lforall[x][(x + \Obj 0) = x]$}
  \RightLabel{\Elim\lforall}
  \UnaryInfC{$(\num{n} + \Obj 0) = \num{n}$}
\end{prooftree}
是$(\num{n} + \Obj 0) = \num{n}$ 的一个自然演绎推导，它有一个未解除假设，而这个未解除假设正是~$\Th{Q}$ 的一条公理。

归纳步骤：设对任意 $n$，已有$\Th{Q} \Proves (\num{n} +
\num{m}) = \num{n+m}$（例如，通过推导 $\delta_{n,m}$）。还需证明$\Th{Q} \Proves (\num{n} + \num{m+1}) =
\num{n+m+1}$。注意到$\num{m+1} \ident \num{m}'$，且$\num{n+m+1} \ident \num{n+m}'$。因此，我们寻求从~$\Th{Q}$ 的公理出发，推导出$(\num{n} + \num{m}') = \num{n+m}'$。我们的推导可以使用归纳假设所保证存在的推导 $\delta_{n,m}$。
\begin{prooftree}
  \AxiomC{}
  \RightLabel{$\delta_{n,m}$}
  \DeduceC{$(\num{n} + \num{m}) = \num{n+m}$}
  \AxiomC{$\lforall[x][\lforall[y][(x+y') = (x+y)']]$}
  \RightLabel{\Elim\lforall}
  \UnaryInfC{$\lforall[y][(\num{n}+y') = (\num{n}+y)']$}
  \RightLabel{\Elim\lforall}
  \UnaryInfC{$(\num{n}+\num{m}') = (\num{n}+\num{m})'$}
    \RightLabel{\Elim=}
    \BinaryInfC{$(\num{n}+\num{m}') = \num{n+m}'$}
\end{prooftree}
在最后的 $\Elim=$ 推理中，我们将右前提右端$(\num{n} + \num{m})'$ 中的子项 $\num{n} + \num{m}$ 替换为项 $\num{n+m}$。
\end{proof}

在 \olref[inc][req][min]{lem:less-zero} 中，我们证明了$\Th{Q} \Proves
\lforall[x][\lnot x < \Obj 0]$。一个实际的推导会是什么样子呢？

\begin{proof}[{\olref[inc][req][min]{lem:less-zero}} 的证明]
为了证明这种全称结论，我们使用 $\Intro\lforall$，它需要先推导出 $\lnot a < \Obj 0$。考察公理 $!Q_8$，这意味着要证明 $\lnot \exists z (z' + a) = \Obj 0$。具体来说，若能证明后者，则由 $!Q_8$ 即可证明前者（回想一下，$A \liff B$ 是 $(A \lif B) \land (B \lif
A)$ 的缩写）。
\begin{prooftree}\footnotesize
  \AxiomC{$\lnot\lexists[z][(z' + a) = \Obj 0]$}
  \AxiomC{$\lforall[x][\lforall[y][(x < y \liff \lexists[z][(z' + x) = y])]]$}
  \RightLabel{\Elim\lforall}
  \UnaryInfC{$\lforall[y][(a < y \liff \lexists[z][(z' + a) = y])]$}
  \RightLabel{\Elim\lforall}
  \UnaryInfC{$a < \Obj 0 \liff \lexists[z][(z' + a) = \Obj 0]$}
  \RightLabel{\Elim\land}
  \UnaryInfC{$a < \Obj 0 \lif \lexists[z][(z' + a) = \Obj 0]$}
  \AxiomC{$\Discharge{a < \Obj 0}{1}$}
  \RightLabel{\Elim\lif}
  \BinaryInfC{$\lexists[z][(z' + a) = \Obj 0]$}
  \RightLabel{\Elim\lnot}
  \insertBetweenHyps{\hskip -3em}
  \BinaryInfC{$\lfalse$}
  \DischargeRule{\Intro\lnot}{1}
  \UnaryInfC{$\lnot a<\Obj 0$}
\end{prooftree}
这是从 $\lnot\lexists[z][(z' +
a) = \Obj 0]$（以及公理~$!Q_8$）出发的 $\lnot a<\Obj 0$ 的推导；我们称它为~$\delta_1$。

现在，我们如何从~$\Th{Q}$ 的公理证明 $\lnot\lexists[z][(z' + a) = \Obj 0]$？要证明像这样的否定式，我们需要如下形式的推导：
\begin{prooftree}
  \AxiomC{$\Discharge{\lexists[z][(z' + a) = \Obj 0]}{2}$}
  \DeduceC{$\lfalse$}
  \DischargeRule{\Intro\lnot}{2}
  \UnaryInfC{$\lnot\lexists[z][(z' + a) = \Obj 0]$}
  \end{prooftree}
为了从一个存在性断言得到矛盾，我们为存在量化的变元~$z$ 引入一个常项~$b$，并使用 \Elim\lexists：
\begin{prooftree}
  \AxiomC{$\Discharge{\lexists[z][(z' + a) = \Obj 0]}{2}$}
  \AxiomC{$\Discharge{(b'+a) = \Obj 0}{3}$}
  \RightLabel{$\delta_2$}
    \DeduceC{$\lfalse$}
    \DischargeRule{\Elim\exists}{3}
    \BinaryInfC{$\lfalse$}
  \DischargeRule{\Intro\lnot}{2}
  \UnaryInfC{$\lnot\lexists[z][(z' + a) = \Obj 0]$}
  \end{prooftree}
现在的任务是填充 $\delta_2$，即从 $(b'+a) = \Obj 0$ 和~$\Th{Q}$ 的公理证明 $\lfalse$。$Q_2$ 说 $\Obj 0$ 不可能是某个数的后继，所以一种方法是证明 $(b' + a)$ 等于某个数的后继。由于这个表达式本身是一个和式，加法的公理必须发挥作用。若 $\eq[a][\Obj 0]$，则由 $Q_5$ 可得 $\eq[(b'
+ a)][b']$，也就是说 $b' + a$ 是某个数（即~$b$）的后继。另一方面，如果对某个 $c$ 有 $\eq[a][c']$，那么由 \Elim\eq 可得 $\eq[(b'+a)][(b'+c')]$，再由~$Q_6$ 可得 $\eq[(b'+c')][(b'+c)']$。因此 $b'+a$ 同样是某个数的后继——这次是 $b'+c$。因此，策略是将任务分为这两种情况。我们还必须验证 $\Th{Q}$ 能证明其中一种情况成立，即 $\Th{Q} \Proves a = 0 \lor \lexists[y][(a = y')]$，但这直接由 $Q_3$ 通过 \Elim\lforall 得出。以下是两种情况：

情况 1：从 $\eq[a][\Obj 0]$ 和 $\eq[(b'+a)][\Obj
  0]$（以及公理 $Q_2$、$Q_5$）证明 $\lfalse$：
\begin{prooftree}\footnotesize
  \AxiomC{$\lforall[x][\lnot \Obj 0 = x']$}
  \RightLabel{\Elim\lforall}
  \UnaryInfC{$\lnot \Obj 0 = b'$}
  \AxiomC{$\lforall[x][(x+\Obj 0) = x]$}
  \RightLabel{\Elim\lforall}
  \UnaryInfC{$(b' + \Obj 0) = b'$}
  \AxiomC{$a = \Obj 0$}
  \AxiomC{$(b'+a) = \Obj 0$}
  \RightLabel{\Elim=}
  \BinaryInfC{$(b' + \Obj 0) = \Obj 0$}
  \doubleLine
  \UnaryInfC{$\Obj 0 = (b' + \Obj 0)$}
      \insertBetweenHyps{\hskip -.5em}
  \RightLabel{\Elim=}
  \BinaryInfC{$\Obj 0 = b'$}
  \RightLabel{\Elim\lnot}
  \BinaryInfC{$\lfalse$}
\end{prooftree}
称这个推导为~$\delta_3$。（我们通过一条双推理线，简写了从 $(b' + \Obj 0) = \Obj 0$ 到 $\Obj 0 = (b' + \Obj 0)$ 的推导。）

情况 2：从 $\lexists[y][a = y']$ 和 $\eq[(b'+a)][\Obj 0]$（以及公理 $Q_2$、$Q_6$）证明 $\lfalse$。我们首先展示如何从 $\eq[a][c']$ 和 $\eq[(b'+a)][\Obj 0]$ 推导出 $\lfalse$。
\begin{prooftree}\footnotesize
  \AxiomC{$\lforall[x][\lnot \Obj 0 = x']$}
  \RightLabel{\Elim\lforall}
  \UnaryInfC{$\lnot \Obj 0 = (b'+c)'$}
  \AxiomC{$a = c'$}
  \AxiomC{$(b'+a) = \Obj 0$}
  \RightLabel{\Elim=}
          \insertBetweenHyps{\hskip -.3em}
  \BinaryInfC{$\Obj (b'+c') = \Obj 0$}
  \AxiomC{$\lforall[x][\lforall[y][(x+y') = (x+y)']]$}
  \RightLabel{\Elim\lforall}
  \UnaryInfC{$\lforall[y][(b'+y') = (b'+y)']$}
  \RightLabel{\Elim\lforall}
  \UnaryInfC{$(b' + c') = (b'+c)'$}
  \RightLabel{\Elim=}
        \insertBetweenHyps{\hskip .5em}
  \BinaryInfC{$\Obj 0 = (b' + c)'$}
  \RightLabel{\Elim\lnot}
          \insertBetweenHyps{\hskip -.5em}
  \BinaryInfC{$\lfalse$}
\end{prooftree}
称此推导为 $\delta_4$。我们通过应用 $\Elim\lexists$ 并解除假设 $\eq[a][c']$，得到所需的推导 $\delta_5$：
\begin{prooftree}
  \AxiomC{$\lexists[y][a=y']$}
  \AxiomC{$\Discharge{a = c'}{6} \quad \eq[(b'+a)][\Obj 0]$}
  \RightLabel{$\delta_4$}
  \DeduceC{$\lfalse$}
  \DischargeRule{\Elim\exists}{6}
  \BinaryInfC{$\lfalse$}
\end{prooftree}

  
将所有部分组合起来，完整的证明如下所示：
\begin{prooftree}\footnotesize
  \AxiomC{$\Discharge{\lexists[z][(z' + a) = \Obj 0]}{2}$}
  \AxiomC{$\begin{gathered}
  \lforall[x][(x = 0 \lor {}]\\
  \lexists[y][(a = y')])
  \end{gathered}$}
  \RightLabel{\Elim\lforall}
  \UnaryInfC{$\begin{gathered}a = 0 \lor {}\\
  \lexists[y][(a = y')]
  \end{gathered}$}
  \AxiomC{$\begin{gathered}[b]
      \Discharge{a = \Obj 0}{7} \\
      \Discharge{(b'+a) = \Obj 0}{3}
      \end{gathered}$}
    \RightLabel{$\delta_3$}
    \DeduceC{$\lfalse$}
    \AxiomC{$\begin{gathered}[b]
        \Discharge{\lexists[y][a=y']}{7} \\
        \Discharge{(b'+a) = \Obj 0}{3}
        \end{gathered}$\quad}
    \RightLabel{$\delta_5$}
    \DeduceC{$\lfalse$}
    \DischargeRule{\Elim\lor}{7}
    \insertBetweenHyps{\hskip -.5em}
    \TrinaryInfC{$\lfalse$}
    \RightSubproofLabel{$\delta_2$}
    \DischargeRule{\Elim\exists}{3}
    \BinaryInfC{$\lfalse$}
  \DischargeRule{\Intro\lnot}{2}
  \UnaryInfC{$\lnot\lexists[z][(z' + a) = \Obj 0]$}
  \RightLabel{$\delta_1$}
  \DeduceC{$\lnot a<\Obj 0$}
  \RightLabel{\Intro\lforall}
  \UnaryInfC{$\lforall[x][\lnot x < \Obj 0]$}
\end{prooftree}
\end{proof}

在 \olref[inc][inp][ros]{thm:rosser} 的证明中，我们将
$\ORProv(y)$ 定义为 \[\lexists[x][(\OPrf(x, y) \land \lforall[z][(z < x
    \lif \lnot \ORefut(z, y))])].\]其中
$\OPrf(x,y)$ 是在 $\Th{Q}$ 中表示理论~$\Th{T}$（一个一致且可公理化的 $\Th{Q}$ 的扩张）的证明关系的公式，而 $\ORefut(z, y)$ 是表示反驳关系的公式。
这意味着，如果 $n$ 是~$!A$ 的一个证明的Gödel数，那么 $\Th{Q} \Proves
\OPrf(\num{n}, \gn{!A})$，否则 $\Th{Q} \Proves
\lnot\OPrf(\num{n}, \gn{!A})$。
类似地，如果 $n$ 是 $\lnot !A$ 的一个证明的Gödel数，那么 $\Th{Q} \Proves \ORefut(\num{n},
\gn{!A})$，否则 $\Th{Q} \Proves \lnot\ORefut(\num{n},
\gn{!A})$。
我们利用对角线引理找到一个语句 $!R$，使得 $\Th{Q} \Proves !R \liff \lnot \ORProv(\gn{!R})$。
Rosser定理断言 $\Th{T} \Proves/ !R$ 且 $\Th{T} \Proves/ \lnot !R$。
这两个结论都是间接证明的：我们证明，若 $\Th{T} \Proves !R$，则 $\Th{T}$ 不一致，即 $\Th{T} \Proves \lfalse$；若 $\Th{T} \Proves \lnot !R$，同样如此。

\begin{proof}[Proof of {\olref[inc][inp][ros]{thm:rosser}}]
首先，我们证明关于~$<$ 的一些事实。由
\olref[inc][req][min]{lem:less-nsucc}，我们知道对每个~$n$，$\Th{Q} \Proves
\lforall[x][(x < \num {n+1} \lif (\eq[x][\Obj 0] \lor \dots \lor
  \eq[x][\num n]))]$。因此当然（若 $n>1$）也有
$\Th{Q} \Proves \lforall[x][(x < \num {n} \lif (\eq[x][\Obj 0] \lor
  \dots \lor \eq[x][\num{n-1}]))]$。我们可以利用这一点从 $a < \num{n}$ 推导出
$\eq[a][\Obj 0] \lor \dots \lor \eq[a][\num{n-1}]$：
\begin{prooftree}
  \AxiomC{$a < \num{n}$}
  \AxiomC{}
  \DeduceC{$\lforall[x][(x < \num{n} \lif (x = \num{0} \lor 
      \dots \lor x = \num{n-1}))]$}
  \RightLabel{\Elim\forall}
  \UnaryInfC{$a < \num{n} \lif (a = \num{0} \lor 
      \dots \lor a = \num{n-1})$}
  \RightLabel{\Elim\lif}
  \BinaryInfC{$a = \num{0} \lor \dots \lor a = \num{n-1}$}
\end{prooftree}
我们称这个推导为 $\lambda_1$。

现在，为证明 $\Th{T} \Proves/ !R$，我们假设 $\Th{T}
\Proves !R$（有一个推导~$\delta$），并证明这将导致 $\Th{T}$ 不一致。令 $n$ 为~$\delta$ 的Gödel数。由于 $\OPrf$ 在~$\Th{Q}$ 中表示证明关系，所以存在 $\OPrf(\num{n},
\gn{!R})$ 的一个推导~$\delta_1$。此外，由于假定 $\Th{T}$ 是一致的，没有 $k < n$ 是~$!R$ 的反驳的Gödel数，因此对每个 $k < n$，都有 $\Th{Q} \Proves \lnot \ORefut(\num{k}, \gn{!R})$；令
$\rho_k$ 为相应的推导。我们得到 $\ORProv(\gn{!R})$ 的一个推导：
\begin{prooftree}\footnotesize
  \AxiomC{}
  \RightLabel{$\delta_1$}
  \DeduceC{$\OPrf(\num{n}, \gn{!R})$}

  \AxiomC{$\Discharge{a < \num{n}}{1}$}
  \RightLabel{$\lambda_1$}
  \DeduceC{$\begin{gathered}[b]
      a= \num{0} \lor \dots {} \\
      {} \lor a = \num{n-1}
      \end{gathered}$}
  \AxiomC{$\dots$}

  \AxiomC{$\Discharge{a=\num{k}}{2}$}
  \AxiomC{}
  \RightLabel{$\rho_k$}
  \DeduceC{$\lnot \ORefut(\num{k}, \gn{!R})$}
  \RightLabel{\Elim=}
  \BinaryInfC{$\lnot \ORefut(a, \gn{!R})$}
  \AxiomC{$\dots$}
  \DischargeRule{$\Elim\lor^*$}{2}
  \doubleLine
  \insertBetweenHyps{\hskip -1pt}
  \QuaternaryInfC{$\lnot \ORefut(a, \gn{!R})$}
  \DischargeRule{\Intro\lif}{1}
  \UnaryInfC{$a < \num{n} \lif \lnot \ORefut(a, \gn{!R})$}
  \RightLabel{\Intro\lforall}
  \UnaryInfC{$\lforall[z][(z < \num{n} \lif \lnot \ORefut(z, \gn{!R}))]$}

    \insertBetweenHyps{\hskip -5pt}
  \RightLabel{\Intro\land}
  \BinaryInfC{$\OPrf(\num{n}, \gn{!R}) \land \lforall[z][(z < \num{n} \lif \lnot \ORefut(\num{z}, \gn{!R}))]$}
  \RightLabel{\Intro\lexists}
  \UnaryInfC{$\lexists[x][(\OPrf(x, \gn{!R}) \land \lforall[z][(z < x \lif \lnot \ORefut(z, \gn{!R}))])]$}
\end{prooftree}
（上面我们用 $\Elim\lor^*$ 来缩写多次应用 $\Elim\lor$。）我们已经证明了，如果 $\Th{T} \Proves !R$，则存在~$\ORProv(\gn{!R})$ 的一个推导。那么，由于 $\Th{T} \Proves R \liff
\lnot \ORProv(\gn{!R})$，也有 $\Th{T} \Proves \ORProv(\gn{!R}) \lif
\lnot R$，我们将有 $\Th{T} \Proves \lnot !R$，从而 $\Th{T}$ 不一致。

现在证明 $\Th{T} \Proves/ \lnot !R$。同样，假设如此。那么就存在 $\lnot !R$ 的一个推导 $\rho$，其Gödel数为 $m$——即~$!R$ 的一个反驳——从而由推导~$\rho_1$ 可得 $\Th{Q} \Proves
\ORefut(\num{m}, \gn{!R})$。由于我们假设 $\Th{T}$ 是一致的，故 $\Th{T} \Proves/ !R$。所以对所有 $k$，$k$ 都不是~$!R$ 的推导的Gödel数，因此由推导~$\pi_k$ 可得 $\Th{Q} \Proves \lnot
\OPrf(\num{k}, \gn{!R})$。于是我们有：
  
\begin{prooftree}\footnotesize
  \AxiomC{}
  \RightLabel{$\lambda_2$}
  \DeduceC{$\begin{gathered}[b] a = \num{0} \lor \dots \lor\\
      a = \num{m} \lor \num{m} < a\end{gathered}$}
  \AxiomC{$\dots$}
  \AxiomC{$\begin{gathered}\Discharge{\OPrf(a, \gn{!R})}{1}\\
    \Discharge{a = \num{k}}{2}\end{gathered}$}
  \RightLabel{$\pi_k'$}
  \DeduceC{$\lfalse$}
  \RightLabel{$\lfalse_I$}
  \UnaryInfC{$\num{m}<a$}
  \AxiomC{$\dots$}
  \AxiomC{$\Discharge{\num{m} < a}{2}$}
  \DischargeRule{$\Elim\lor^*$}{2}
  \insertBetweenHyps{\hspace{-.1em}}
  \doubleLine
  \QuinaryInfC{$\num{m} < a$}
  \AxiomC{}
  \RightLabel{$\rho_1$}
  \DeduceC{$\ORefut(\num{m}, \gn{!R})$}
  \RightLabel{\Intro\land}
    \insertBetweenHyps{\hspace{-.1em}}
  \BinaryInfC{$\num{m}
    < a \land \ORefut(\num{m}, \gn{!R})$}
  \RightLabel{\Intro\exists}
  \UnaryInfC{$\exists z(z
    < a \land \ORefut(z, \gn{!R}))$}
  \DischargeRule{\Intro\lif}{1}
  \UnaryInfC{$\OPrf(a, \gn{!R}) \lif \exists z(z
    < a \land \ORefut(z, \gn{!R}))$}
  \RightLabel{\Intro\forall}
  \UnaryInfC{$\forall x(\OPrf(x, \gn{!R}) \lif \exists z(z
    < x \land \ORefut(z, \gn{!R})))$}
  \DeduceC{$\lnot\exists x(\OPrf(x, \gn{!R}) \land \forall z(z
    < x \lif \lnot \ORefut(z, \gn{!R})))$}
\end{prooftree}
其中 $\pi_k'$ 是如下推导
\begin{prooftree}
  \AxiomC{}
  \RightLabel{$\pi_k$}
  \DeduceC{$\lnot \OPrf(\num{k}, \gn{!R})$}
  \AxiomC{$a = \num{k}$}
  \AxiomC{$\OPrf(a, \gn{!R})$}
  \RightLabel{\Elim=}
  \BinaryInfC{$\OPrf(\num{k}, \gn{!R})$}
  \RightLabel{\Elim\lnot}
  \BinaryInfC{$\lfalse$}
\end{prooftree}
而 $\lambda_2$ 是
\begin{prooftree}\footnotesize
  \AxiomC{}
  \RightLabel{$\lambda_3$}
  \DeduceC{$\begin{gathered}[b](a < \num{m} \lor {}\\a = \num{m}) \lor {}\\ \num{m} < a\end{gathered}$}

  \AxiomC{$\Discharge{a < \num{m}}{3}$}
  \RightLabel{$\lambda_1$}
  \DeduceC{$\begin{gathered}[b]a = \num{0} \lor \dots \lor {}\\
  a = \num{m-1}\end{gathered}$}
  %\RightLabel{\Intro\lor}
  \doubleLine
  \UnaryInfC{$\begin{gathered}[b]
      a = \num{0} \lor \dots \lor \\
      a = \num{m} \lor \num{m} < a
    \end{gathered}$}

  \AxiomC{$\Discharge{a = \num{m}}{3}$}
  %\RightLabel{\Intro\lor}
  \doubleLine
  \UnaryInfC{$\begin{gathered}[b]
      a = \num{0} \lor \dots \lor \\
      a = \num{m} \lor \num{m} < a
    \end{gathered}$}
  
  \AxiomC{$\Discharge{\num{m} < a}{3}$}
  \doubleLine
  \RightLabel{$\Intro\lor^*$}
  \UnaryInfC{$\begin{gathered}[b]
      a = \num{0} \lor \dots \lor \\
      a = \num{m} \lor \num{m} < a
    \end{gathered}$}

  %\insertBetweenHyps{\hskip 5pt}
  \doubleLine
  \DischargeRule{$\Elim\lor^2$}{3}
  \QuaternaryInfC{$a = \num{0} \lor \dots \lor a = \num{m} \lor \num{m} < a$}
\end{prooftree}
（推导 $\lambda_3$ 的存在性由 \olref[inc][req][min]{lem:trichotomy} 保证。我们用 $\Intro\lor^*$ 缩写重复使用 $\Intro\lor$，并将从 $(a < \num{m} \lor a = \num{m}) \lor \num{m} < a$ 推出 $a = \num{0} \lor \dots \lor a = \num{m} \lor \num{m} < a$ 时双重使用 $\Elim\lor$ 记为 $\Elim\lor^2$。）
\end{proof}

\OLEndChapterHook
